\chapter{智能体与工具使用}

\section{从对话到行动}

在前面的章节中，我们学习了提示词工程和 RAG 技术。这些技术让 LLM 能够更好地理解和回答问题。但是，LLM 本身只能生成文本，无法直接与外部世界交互。

想象你是一个被关在房间里的人，你知识渊博，但你无法打电话、无法上网、无法操作电脑。你只能通过门缝递纸条与人交流。这就是 LLM 的处境——它有强大的推理能力，但缺乏执行能力。

智能体（Agent）技术就是给 LLM 装上"手和脚"，让它不仅能思考，还能行动。

\section{什么是智能体}

智能体（Agent）是一个能够感知环境并采取行动以实现目标的系统。在 AI 领域，LLM 智能体通常包含以下组件：

\begin{itemize}
	\item \textbf{大脑}：LLM 作为核心推理引擎
	\item \textbf{工具}：可以调用的外部功能（如搜索引擎、代码执行器、API等）
	\item \textbf{记忆}：短期记忆（对话历史）和长期记忆（文件系统）
	\item \textbf{规划}：将复杂任务分解为可执行的步骤
\end{itemize}

打个比方，LLM 是一个聪明的大脑，智能体就是给这个大脑配上工具箱、记事本和计划表，让它能够完成实际任务。

\section{Agent 循环：智能体的心跳}

\subsection{ReAct 框架}

ReAct（Reasoning and Acting）是最经典的智能体框架，它将推理和行动结合起来。
ReAct 的核心思想是让 LLM 交替进行"思考"（Reasoning）和"行动"（Acting），
形成一个 Thought-Action-Observation 的循环：

\begin{verbatim}
用户：帮我查一下特斯拉当前的股价，并分析最近走势

Thought 1: 我需要先获取特斯拉的当前股价信息
Action 1: 调用股票查询工具，查询 TSLA
Observation 1: 特斯拉当前股价 245.50 美元

Thought 2: 现在我需要获取最近的历史数据来分析走势
Action 2: 调用历史数据工具，查询 TSLA 近30天数据
Observation 2: [返回30天的股价数据]

Thought 3: 根据数据，我可以分析走势了
Final Answer: 特斯拉当前股价 245.50 美元，
最近30天呈现上涨趋势，涨幅约 11.6%...
\end{verbatim}

\subsection{真实的 Agent 循环}

上面的 ReAct 示例做了简化。在实际的智能体系统中（如 Claude Code），Agent 循环要复杂得多。
下面是一个真实 Agent 循环的核心流程：

\begin{enumerate}
	\item \textbf{构建系统提示词}：组装系统提示词，包含工具定义、环境信息、记忆等
	\item \textbf{调用模型}：将对话历史发送给 LLM，获取模型响应
	\item \textbf{解析工具调用}：检查模型输出中是否包含工具调用请求
	\item \textbf{权限检查}：在执行工具前，验证用户是否授权
	\item \textbf{执行工具}：调用相应的工具函数，获取结果
	\item \textbf{注入结果}：将工具执行结果作为新消息加入对话历史
	\item \textbf{判断终止}：如果没有工具调用，说明模型认为任务完成，退出循环；否则回到第2步
\end{enumerate}

这个过程可以用下面的伪代码表示：

\begin{verbatim}
function agentLoop(messages, tools, systemPrompt):
    while true:
        // 1. 调用 LLM
        response = callModel(messages, tools, systemPrompt)

        // 2. 检查是否有工具调用
        if response has no tool_use blocks:
            return response  // 任务完成，退出循环

        // 3. 执行所有工具调用
        for each tool_use in response:
            tool = findTool(tool_use.name)
            result = tool.execute(tool_use.input)
            messages.append(tool_result)

        // 4. 继续循环，让模型看到工具结果
\end{verbatim}

这个循环看起来简单，但实际实现中有许多精妙的设计。

\subsection{流式工具执行}

在 Claude Code 中，工具执行并不是等模型完全输出后才开始的。
当模型的输出以流（stream）的形式返回时，一旦某个工具调用的 JSON 参数完整接收，
系统就会立即开始执行该工具，而不需要等待模型输出完毕。

打个比方，这就像餐厅里服务员还没记完你的所有菜，就已经把先点好的菜传给厨房了。
这种\textbf{流式工具执行}（Streaming Tool Execution）大大减少了等待时间。

此外，标记为"并发安全"的工具可以并行执行。比如同时搜索多个文件，
不需要等一个搜索完了再搜索下一个。

\subsection{错误恢复}

真实的 Agent 循环需要处理各种异常情况。Claude Code 中实现了多层恢复机制：

\begin{itemize}
	\item \textbf{输出截断}：当模型输出达到 token 上限时，自动提高限制并重试（最多3次）
	\item \textbf{上下文过长}：当提示词超过模型上下文窗口时，自动压缩对话历史后重试
	\item \textbf{API 错误}：网络或服务端错误时，自动切换到备用模型重试
\end{itemize}

这些恢复机制对用户是透明的，确保智能体能够稳定运行。

\subsection{步数限制与"末日循环"检测}

一个容易被忽视的问题是：如果智能体陷入死循环怎么办？
比如反复调用同一个工具、反复修改同一个文件，却永远无法完成任务。

OpenCode 的解决方案是设置\textbf{步数上限}（默认25步）。
当智能体连续执行超过这个上限时，系统会强制终止循环，
并向用户报告当前进度。这就像给智能体设了一个"最大工作时间"，
超时就必须停下来汇报。

Gemini CLI 则实现了\textbf{"末日循环"检测}（Doom Loop Detection）：
系统会监控智能体的行为模式，如果发现它连续多次执行相同的操作
（比如连续3次调用同一个工具、传入相同的参数），
就会主动介入并提醒用户。

\subsection{持久化输入与会话恢复}

在 OpenCode 的 V2 架构中，引入了一个精巧的设计——\textbf{持久化输入}（Durable Input）。
用户发送的消息不会直接交给模型，而是先存入一个数据库"收件箱"，
然后在安全的边界点被"提升"到对话中。

这就像邮件系统：你发出的邮件先进入对方的收件箱，
对方在方便的时候才打开阅读。
这种设计的好处是，即使进程崩溃，未处理的消息也不会丢失，
重启后可以继续处理。

\section{工具系统}

工具使用（Tool Use）是智能体的核心能力。它允许 LLM 在对话过程中调用外部工具来获取信息或执行操作。

\subsection{工具的定义}

每个工具都需要告诉 LLM 三件事：我叫什么、我能做什么、你需要给我什么参数。
这些信息通过 JSON Schema 来定义：

\begin{verbatim}
{
  "name": "get_weather",
  "description": "获取指定城市的天气信息",
  "input_schema": {
    "type": "object",
    "properties": {
      "city": {
        "type": "string",
        "description": "城市名称"
      },
      "date": {
        "type": "string",
        "description": "日期，如 'today', 'tomorrow'"
      }
    },
    "required": ["city"]
  }
}
\end{verbatim}

在 Claude Code 中，工具的定义更加丰富。一个编码智能体通常配备 60 多个工具，
涵盖了文件操作、代码搜索、命令执行、网络访问等方方面面。
以下是 Claude Code 中的部分核心工具：

\begin{table}[H]
	\centering
	\begin{tabular}{|l|l|}
		\hline
		\textbf{工具名称} & \textbf{功能描述} \\
		\hline
		Bash          & 执行终端命令        \\
		Read          & 读取文件内容        \\
		Write         & 创建或覆盖文件       \\
		Edit          & 精确替换文件中的文本    \\
		Glob          & 按模式搜索文件名      \\
		Grep          & 按内容搜索文件       \\
		Agent         & 启动子智能体        \\
		WebFetch      & 获取网页内容        \\
		WebSearch     & 搜索引擎搜索        \\
		LSP           & 代码智能（跳转定义等）   \\
		\hline
	\end{tabular}
	\caption{Claude Code 核心工具列表}
\end{table}

\subsection{工具调用的完整生命周期}

一个工具调用从开始到结束，要经过以下步骤：

\begin{enumerate}
	\item \textbf{模型决策}：LLM 在输出中生成一个 \texttt{tool\_use} 块，包含工具名和参数
	\item \textbf{工具查找}：系统根据名称在已注册的工具列表中查找对应工具
	\item \textbf{输入验证}：用 JSON Schema 验证模型传入的参数是否合法
	\item \textbf{权限检查}：检查用户是否授权该操作（通过 Hook 系统或用户确认）
	\item \textbf{执行工具}：调用工具的 \texttt{execute} 方法，获取结果
	\item \textbf{结果格式化}：将结果转换为 \texttt{tool\_result} 格式
	\item \textbf{注入对话}：将结果作为用户消息加入对话历史，供模型下次读取
\end{enumerate}

用一个具体的例子来说明。当用户说"帮我看看 main.py 的内容"时：

\begin{verbatim}
// 第1步：模型输出 tool_use 块
{
  "type": "tool_use",
  "id": "toolu_01ABC",
  "name": "Read",
  "input": { "file_path": "/home/user/main.py" }
}

// 第2-5步：系统查找 Read 工具，验证参数，执行读取
// 读取结果：
"def main():
    print('Hello, World!')

if __name__ == '__main__':
    main()"

// 第6-7步：格式化并注入对话历史
{
  "type": "tool_result",
  "tool_use_id": "toolu_01ABC",
  "content": "def main():\n    print('Hello, World!')\n..."
}
\end{verbatim}

模型看到这个结果后，就可以基于文件内容回答用户的问题了。

\subsection{工具执行策略}

不同的智能体系统在工具执行策略上有各自的设计。

\paragraph{并发控制}
在 Claude Code 中，标记为"并发安全"的工具可以并行执行。
而 Gemini CLI 采用了一种更灵活的方式——它在每个工具的参数中自动添加一个
\texttt{wait\_for\_previous} 参数，让模型自己决定哪些工具可以并行、哪些必须串行。

\begin{verbatim}
// Gemini CLI 自动为工具添加的参数
{
  "name": "read_file",
  "parameters": {
    "path": {"type": "string"},
    "wait_for_previous": {
      "type": "boolean",
      "description": "是否等待前一个工具执行完毕"
    }
  }
}
\end{verbatim}

这种设计的巧妙之处在于：把并发控制的决策权交给了模型。
模型可以根据任务依赖关系，智能地决定执行顺序。

\paragraph{后台执行}
Gemini CLI 的 Shell 工具支持后台执行。当命令耗时较长时（比如编译大型项目），
智能体可以让命令在后台运行，同时继续处理其他任务。
系统通过 PID 追踪来管理后台进程，确保能够获取最终结果。

\begin{verbatim}
// 后台执行示例
{
  "command": "make build",
  "is_background": true
}
// 系统返回：进程已在后台启动，PID: 12345
\end{verbatim}

\paragraph{工具输出管理}
工具的输出可能非常长（比如读取一个大文件）。OpenCode 实现了\textbf{有界输出}机制：
工具结果被分为"预览"和"完整内容"两部分。预览保留在对话历史中供模型参考，
完整内容则写入临时文件。如果模型需要完整内容，可以通过专门的工具读取。

这就像看报告时，摘要直接附在邮件正文里，完整报告作为附件。
收件人先看摘要，需要时再打开附件。

\subsection{权限与安全}

智能体拥有执行能力后，安全性就变得至关重要。
不同的智能体系统实现了不同层次的权限控制。

\paragraph{审批模式}
Gemini CLI 定义了四种审批模式，从严格到宽松：

\begin{itemize}
	\item \textbf{DEFAULT}：每次工具调用都需要用户确认
	\item \textbf{AUTO\_EDIT}：自动批准文件编辑操作，其他操作仍需确认
	\item \textbf{PLAN}：计划模式，智能体只能读取和分析，不能修改文件
	\item \textbf{YOLO}：自动批准所有操作（适合完全信任的场景）
\end{itemize}

这种分级设计让用户可以根据任务的危险程度选择合适的模式。
比如日常编码可以用 AUTO\_EDIT，而处理敏感数据时用 DEFAULT。

\paragraph{IAM 风格的权限规则}
OpenCode 采用了类似云服务的 IAM（Identity and Access Management）权限模型。
每个智能体都有自己的权限规则集：

\begin{verbatim}
// OpenCode 的权限规则示例
{
  "effect": "deny",
  "action": "bash",
  "resource": "rm -rf *"
}
\end{verbatim}

这种规则可以在智能体启动时就过滤掉不允许的工具，
而不是等到执行时才发现权限不足。
这就像门禁系统：不给你某间房的门禁卡，你就根本进不去，
而不是进去了才发现门被锁了。

\paragraph{平台特定的沙箱}
Gemini CLI 实现了平台特定的沙箱系统，针对 Linux、macOS、Windows
分别实现了不同的安全隔离机制。这包括文件系统访问限制、
网络访问控制、进程隔离等。

\paragraph{主动权限建议}
一个有趣的创新是 Gemini CLI 的\textbf{主动权限建议}功能。
当智能体检测到某个操作可能需要特定权限时，它会主动建议用户授权，
而不是等到操作失败后才提示。这提升了用户体验，减少了中断。

\section{函数调用：连接模型与工具的桥梁}

\subsection{函数调用的工作原理}

函数调用（Function Calling）是工具使用的底层实现机制。
我们在第9章学习了 Transformer 的解码过程——模型逐词生成文本。
函数调用在此基础上做了扩展：模型的输出不再仅仅是文本，还可以是结构化的工具调用指令。

具体来说，当我们将工具定义发送给 API 时，模型会在词表中增加一些特殊的 token，
如 \texttt{tool\_use} 标记。当模型决定调用工具时，它会生成这些特殊 token，
而不是普通的文字。

\begin{verbatim}
// 发送给 API 的请求
{
  "model": "claude-sonnet-4-20250514",
  "messages": [{"role": "user", "content": "计算 123 * 456"}],
  "tools": [{
    "name": "calculate",
    "description": "执行数学计算",
    "input_schema": {
      "type": "object",
      "properties": {
        "expression": {"type": "string"}
      },
      "required": ["expression"]
    }
  }]
}

// API 返回的响应
{
  "content": [
    {"type": "text", "text": "让我帮你计算一下。"},
    {"type": "tool_use", "id": "toolu_01ABC",
     "name": "calculate",
     "input": {"expression": "123 * 456"}}
  ]
}
\end{verbatim}

\subsection{多轮工具调用}

一个强大的智能体往往需要连续调用多个工具。每次工具调用的结果会被加入对话历史，
模型基于更新后的历史继续推理，形成多轮交互：

\begin{verbatim}
用户：帮我创建一个 Python 项目

第1轮 - 模型调用 Bash 工具：mkdir my_project
第1轮 - 工具返回：创建成功

第2轮 - 模型调用 Write 工具：写入 main.py
第2轮 - 工具返回：写入成功

第3轮 - 模型调用 Write 工具：写入 requirements.txt
第3轮 - 工具返回：写入成功

第4轮 - 模型调用 Bash 工具：python main.py 验证
第4轮 - 工具返回：Hello, World!

第5轮 - 模型输出文本：项目已创建完成，包含以下文件...
（没有工具调用，循环结束）
\end{verbatim}

\section{系统提示词：智能体的灵魂}

\subsection{系统提示词的作用}

如果说工具是智能体的"手和脚"，那么系统提示词（System Prompt）就是智能体的"灵魂"。
它定义了智能体的身份、行为准则和工作方式。

系统提示词与普通提示词的区别在于：它在每次对话开始时就设定好了，
贯穿整个对话过程，模型会始终遵循其中的指示。

\subsection{系统提示词的构建}

不同的智能体系统在系统提示词的构建上采用了不同的策略。

\paragraph{Claude Code：模块化组装}
在 Claude Code 中，系统提示词由多个"片段"动态组装而成：

\begin{enumerate}
	\item \textbf{介绍部分}：定义模型的身份
	\item \textbf{行为准则}：规定模型应该如何行事
	\item \textbf{工具说明}：描述每个工具的使用方法
	\item \textbf{环境信息}：当前日期、工作目录、操作系统等
	\item \textbf{项目上下文}：CLAUDE.md 中的项目约定
	\item \textbf{记忆内容}：智能体记住的用户偏好和项目知识
\end{enumerate}

\paragraph{Gemini CLI：分层记忆系统}
Gemini CLI 实现了\textbf{分层记忆}（Hierarchical Memory）：

\begin{itemize}
	\item \textbf{全局记忆}：\texttt{\~/.gemini/GEMINI.md}，适用于所有项目
	\item \textbf{扩展记忆}：来自 VSCode 等 IDE 的记忆
	\item \textbf{项目记忆}：\texttt{.gemini/GEMINI.md}，仅适用于当前项目
\end{itemize}

这种分层设计让智能体能够区分通用知识和项目特定知识。

\paragraph{OpenCode：Provider 特定的提示词}
OpenCode 支持多个 LLM 提供商（Anthropic、OpenAI、Gemini 等），
不同模型的"性格"和能力不同，因此 OpenCode 为每个 Provider 准备了不同的系统提示词模板：

\begin{verbatim}
function provider(model) {
  if (model.includes("gpt-4")) return PROMPT_GPT4
  if (model.includes("gemini")) return PROMPT_GEMINI
  if (model.includes("claude")) return PROMPT_CLAUDE
  return PROMPT_DEFAULT
}
\end{verbatim}

这种设计的理由是：不同模型对指令的理解方式不同，
针对性的提示词能更好地发挥每个模型的优势。

\paragraph{上下文代数}
OpenCode 的 V2 架构引入了\textbf{系统上下文代数}（System Context Algebra）。
系统提示词不再是一个字符串，而是由多个"上下文源"（Context Source）组合而成的代数结构。
每个上下文源可以独立更新，系统会自动计算最终的提示词。

这就像化学反应：不同的"元素"（上下文源）按照一定的"化学方程式"（组合规则）
反应生成最终的"化合物"（系统提示词）。当某个元素变化时，只需要重新计算，
而不需要从头开始。

\section{子智能体：分工协作}

\subsection{为什么需要子智能体}

复杂的任务往往超出单个智能体的能力范围。就像一个公司需要不同部门协作一样，
智能体也需要"分身"来并行处理不同的子任务。

例如，当用户说"帮我重构认证模块并更新测试"时，主智能体可以：
\begin{itemize}
	\item 启动一个子智能体去分析认证模块的代码结构
	\item 同时启动另一个子智能体去查找相关的测试文件
	\item 收集两个子智能体的结果后，再制定重构方案
\end{itemize}

\subsection{子智能体的隔离}

子智能体在独立的上下文中运行，这就像一个员工被派到一间独立的办公室工作——
他有自己的工作台（消息历史），看不到其他员工的工作，
但他继承了公司的规章制度（系统提示词）。

在 Claude Code 中，子智能体的上下文隔离通过以下机制实现：

\begin{itemize}
	\item \textbf{独立的消息历史}：子智能体有自己的对话记录，不会污染主对话
	\item \textbf{独立的文件缓存}：子智能体读取文件时建立自己的缓存
	\item \textbf{关联的中止控制}：主智能体中止时，子智能体也会自动中止
	\item \textbf{独立的 Agent ID}：每个子智能体有唯一标识，便于追踪
\end{itemize}

\subsection{多智能体协调模式}

在更复杂的场景下，Claude Code 支持"协调者-工人"（Coordinator-Worker）模式：

\begin{itemize}
	\item \textbf{协调者}（Coordinator）：负责理解任务、分配工作、综合结果。
	      它不直接执行具体操作，而是通过派遣工人来完成任务
	\item \textbf{工人}（Worker）：接收具体任务，拥有完整的工具集，
	      独立执行任务后将结果报告给协调者
\end{itemize}

这就像一个项目经理和开发团队的关系：项目经理负责需求分析和任务分配，
开发人员负责具体实现，完成后向项目经理汇报。

\subsection{Agent 间通信协议}

当多个智能体需要协作时，它们之间如何通信是一个关键问题。
Google 提出了 \textbf{A2A}（Agent-to-Agent）协议，用于规范智能体之间的通信。

在 Gemini CLI 中，子智能体有一个重要的限制：\textbf{不允许递归调用}。
也就是说，子智能体不能再启动子智能体。这是为了防止无限嵌套导致的资源失控。

\begin{verbatim}
主智能体
├── 子智能体 A（分析代码结构）
│   └── ❌ 不能再启动子智能体
├── 子智能体 B（查找测试文件）
│   └── ❌ 不能再启动子智能体
└── 汇总结果
\end{verbatim}

这种设计虽然简单，但有效地避免了"俄罗斯套娃"式的无限嵌套问题。

\section{上下文管理}

上下文管理是智能体系统中一个容易被低估的难题。
它不仅关乎对话历史的存储，更影响智能体的推理质量和成本。

\subsection{上下文压缩}

当对话历史接近模型的上下文窗口上限时，所有智能体系统都会进行压缩。
但压缩的策略各有不同。

Claude Code 的压缩相对直接：将早期消息总结为摘要，保留最近的对话。

OpenCode 的压缩更加结构化，它会生成一个包含以下字段的摘要：
\begin{itemize}
	\item \textbf{目标}：用户想要完成什么
	\item \textbf{约束}：有哪些限制条件
	\item \textbf{进展}：已完成、进行中、被阻塞的任务
	\item \textbf{关键决策}：做出了哪些重要决定
	\item \textbf{下一步}：接下来应该做什么
	\item \textbf{相关文件}：涉及哪些文件
\end{itemize}

这种结构化的压缩能更好地保留关键信息，避免智能体在压缩后"忘记"重要上下文。

\subsection{历史加固}

Gemini CLI 实现了一个独特的功能——\textbf{历史加固}（History Hardening）。
在将对话历史发送给模型之前，系统会对历史进行"清洗"，
移除或修正可能误导模型的内容。

打个比方，这就像在把会议纪要发给参会者之前，
先由助理检查一遍，修正其中的错别字和歧义表述，
确保每个人看到的都是清晰准确的信息。

\subsection{上下文纪元}

OpenCode V2 引入了\textbf{上下文纪元}（Context Epoch）的概念。
一个纪元是指系统提示词保持不变的一段时间跨度。
当发生以下事件时，会开启新的纪元：
\begin{itemize}
	\item 上下文压缩
	\item 切换模型或提供商
	\item 切换智能体
\end{itemize}

每个纪元都有一个"基线"（Baseline），记录了该纪元开始时的系统上下文状态。
这就像版本控制系统中的"快照"——你可以随时回到某个快照的状态。

这种设计的好处是：在一个纪元内，系统提示词是不可变的，
这保证了模型推理的一致性。当需要变更时，通过纪元切换来管理，
而不是在对话中途偷偷修改提示词。

\subsection{图结构上下文}

Gemini CLI 采用了最复杂的上下文管理方案——\textbf{图结构上下文}（Graph-based Context）。
传统的上下文管理是一个线性的消息列表，而 Gemini CLI 将上下文表示为一个有向图：

\begin{itemize}
	\item \textbf{节点}（Node）：每条消息是一个节点
	\item \textbf{处理器}（Processor）：对节点进行变换的管道
	\item \textbf{工作缓冲区}（Working Buffer）：处理器可以修改的可变状态
\end{itemize}

这种设计允许对上下文进行复杂的变换操作，比如：
重新排序消息、注入额外的上下文、过滤不相关的内容等。

打个比方，线性上下文就像一条流水线，原料进去，产品出来。
而图结构上下文更像一个化工厂，原料可以在不同的反应罐之间流转，
经过各种处理，最终生成产品。

\section{记忆系统}

\subsection{短期记忆：对话历史}

短期记忆就是当前对话的消息列表。它让智能体能够理解对话的连贯性。

\begin{verbatim}
用户：帮我查一下北京的天气
智能体：[调用天气工具] 北京今天晴，22°C
用户：那上海呢？  // 智能体从上下文知道还是在问天气
智能体：[调用天气工具] 上海今天多云，25°C
\end{verbatim}

短期记忆面临的一个挑战是\textbf{上下文窗口限制}。
模型的输入长度是有上限的，当对话很长时，早期的内容会被"遗忘"。
为了解决这个问题，Claude Code 实现了\textbf{自动压缩}机制：
当对话历史接近上下文窗口上限时，系统会自动将早期的消息压缩成摘要，
保留关键信息，释放空间给新的内容。

\subsection{长期记忆：文件系统}

长期记忆让智能体能够跨会话记住信息。在 Claude Code 中，长期记忆通过文件系统实现。

具体来说，智能体会在项目目录下维护一个 \texttt{memory/} 文件夹，
其中 \texttt{MEMORY.md} 是记忆的索引文件。智能体在工作过程中学到的重要信息
（如用户偏好、项目约定、常见问题解决方案）会被写入这些文件。

下次启动时，智能体会自动读取这些记忆文件，恢复之前的"记忆"。

\begin{verbatim}
// 记忆文件示例 (~/.claude/projects/xxx/memory/MEMORY.md)
# 项目记忆

## 用户偏好
- 使用 TypeScript 而非 JavaScript
- 偏好函数式编程风格
- 测试框架使用 vitest

## 项目约定
- API 路由放在 src/routes/ 目录
- 数据库操作使用 Prisma ORM
- 提交信息遵循 Conventional Commits
\end{verbatim}

这种设计非常巧妙：它不需要修改模型参数，仅仅通过文件系统就实现了持久化记忆。
就像一个工程师把自己的工作笔记放在抽屉里，下次上班时翻开笔记本就能回忆起之前的工作。

\section{编码智能体：Claude Code}

\subsection{编码智能体的特点}

编码智能体（Coding Agent）是专门用于编程任务的智能体。
相比通用智能体，编码智能体需要更强大的代码理解和操作能力。

Claude Code 是 Anthropic 推出的编码智能体。
它基于 Claude 模型，配备了丰富的开发工具，能够理解整个代码库的上下文，
自主完成代码编写、调试、重构等复杂任务。

\subsection{一个完整的工作流程}

让我们通过一个完整的例子来理解 Claude Code 的工作流程。
假设用户说："帮我修复登录页面的 bug，点击按钮没有反应"。

Claude Code 的执行过程如下：

\begin{verbatim}
[思考] 用户报告登录按钮点击无反应，我需要先找到相关代码
[行动] 调用 Grep 工具，搜索 "login" 相关文件
[观察] 找到 src/pages/Login.tsx

[思考] 需要查看登录页面的代码
[行动] 调用 Read 工具，读取 Login.tsx
[观察] 发现 handleSubmit 函数中有个拼写错误

[思考] 需要确认这个函数的调用方式
[行动] 调用 LSP 工具，查找 handleSubmit 的引用
[观察] 确认是按钮的 onClick 事件处理函数

[思考] 问题已定位，现在修复这个拼写错误
[行动] 调用 Edit 工具，将 "fromData" 改为 "formData"
[观察] 修改成功

[思考] 需要验证修复是否正确
[行动] 调用 Bash 工具，运行测试
[观察] 所有测试通过

[最终回答] 问题已修复。原因是 handleSubmit 函数中
"formData" 被拼写为 "fromData"，导致获取不到表单数据。
已修改 src/pages/Login.tsx 第42行，测试全部通过。
\end{verbatim}

在这个过程中，Claude Code 自主完成了"搜索-阅读-分析-修复-验证"的完整流程，
总共调用了5次工具，经历了5轮思考-行动循环。

\subsection{技能系统}

Claude Code 还实现了\textbf{技能}（Skills）系统，可以看作是针对特定任务的"专家模式"。
每个技能定义了一组专用的提示词和工具集，当触发某个技能时，智能体会切换到该技能的
工作方式。

例如，一个代码审查技能可能包含：
\begin{itemize}
	\item 专门的审查提示词（关注安全漏洞、性能问题等）
	\item 限制只使用只读工具（Read、Grep、Glob），不允许修改文件
	\item 输出格式要求（按严重程度分类列出问题）
\end{itemize}

技能可以通过斜杠命令（如 \texttt{/review}）手动触发，也可以由智能体根据上下文自动选择。

\section{智能体的挑战与展望}

\subsection{当前挑战}

尽管智能体技术发展迅速，但仍面临一些挑战：

\begin{itemize}
	\item \textbf{可靠性}：智能体可能会做出错误的决策或调用错误的工具，
	      特别是在复杂的多步任务中，一个错误可能导致后续步骤全部失败
	\item \textbf{安全性}：需要防止智能体执行危险操作，如删除重要文件、
	      泄露敏感信息等
	\item \textbf{效率}：多步推理和工具调用会增加响应时间和 API 成本
	\item \textbf{可控性}：用户希望了解智能体在做什么，并能在必要时介入
\end{itemize}

\subsection{未来展望}

智能体技术正在快速发展，未来的趋势包括：

\begin{itemize}
	\item \textbf{多智能体协作}：多个专业化的智能体协同工作，
	      就像一个高效的开发团队
	\item \textbf{自主学习}：智能体能够从经验中学习，不断改进策略
	\item \textbf{更丰富的工具生态}：MCP（Model Context Protocol）等标准
	      让工具的接入更加规范化和便捷
	\item \textbf{个性化}：智能体能够适应不同用户的偏好和工作方式，
	      通过记忆系统越用越懂你
\end{itemize}

智能体代表了 AI 从"能说"到"能做"的重要跨越。
从提示词工程解决"怎么问"，到 RAG 解决"知道什么"，
到工具使用解决"能做什么"，再到智能体将这些能力整合在一起——
我们正在见证 AI 从"对话助手"进化为"行动执行者"的历史性转变。
